\documentclass[10pt,a4paper]{amsart}
\usepackage[margin=19mm]{geometry}
\usepackage{amsmath,amssymb,mathtools}
\usepackage[scr=rsfs,cal=euler]{mathalfa}
\usepackage{xfrac}
\usepackage[unicode,hidelinks,bookmarks=false]{hyperref}
\newcommand{\ie}{i.e.\ }
\newcommand{\cf}{cf.\ }
\newcommand{\resp}{respectively\ }
\newcommand{\Subsection}{\subsection}
\providecommand{\nobreakdash}{\nobreak\hbox{-}}
\setlength{\emergencystretch}{2em}
\multlinegap0pt
\usepackage{amssymb}
\usepackage{enumerate}
\usepackage{xcolor}
\newtheorem{prop}{Proposition}
\newtheorem{lem}{Lemma}
\newtheorem{thm}{Theorem}
\usepackage{cleveref}
\crefname{prop}{Proposition}{Propositions}
\crefname{lem}{Lemma}{Lemmas}
\crefname{thm}{Theorem}{Theorems}
\newcommand\Ad{\operatorname{\mathsf{Ad}}}
\newcommand{\G}{\mathsf{G}}
\newcommand{\Hyp}{\mathsf{Hyp}}
\newcommand{\I}{\mathsf{I}}
\newcommand\N{\operatorname{N}}
\newcommand{\can}{\operatorname{\mathsf{can}}}
\renewcommand{\deg}{\operatorname{\mathsf{deg}}}
\newcommand{\disc}{\mathsf{disc}}
\newcommand\ind{\operatorname{\mathsf{ind}}}
\newcommand{\la}{\langle}
\renewcommand{\leq}{\leqslant}
\newcommand{\lla}{\la\!\la}
\newcommand{\mg}[1]{{#1}^{\times}}
\renewcommand{\mod}{\,\mathsf{mod}\,}
\newcommand{\ra}{\rangle}
\newcommand{\rra}{\ra\!\ra}      
\newcommand{\s}{\sigma}
\newcommand{\sq}[1]{{#1}^{\times 2}}
\newcommand{\vf}{\varphi}
\title{$R$-triviality for adjoint classical groups of type $C$}
\author{M. Archita}
\date{Source: arXiv:2602.23072v1 (CC BY 4.0)}
\begin{document}
\maketitle

\noindent\emph{Source and licence.} $R$-triviality for adjoint classical groups of type $C$. Version arXiv:2602.23072v1 (CC BY 4.0), distributed by arXiv under the Creative Commons Attribution 4.0 International licence, http://creativecommons.org/licenses/by/4.0/, as recorded in the arXiv licence metadata for this exact version and shown on the abstract page https://arxiv.org/abs/2602.23072v1. Full licence notice: \texttt{licenses/arxiv-2602.23072-CC-BY-4.0.txt}.\par\smallskip

\noindent\textit{Editorial scope.} Counted unit: the article's thm 6 (source0.tex lines 363--365). The article's complete original proof of it is delivered with it (source0.tex lines 366--386, one \texttt{proof} environment of 21 lines). No mathematics is written, repaired or re-derived: the statement and the proof are the article's own lines, in the article's own notation, and no retained line is substituted or re-flowed.  Retained uncounted as the source-local context the proof needs: lines 283--286; lines 340--343.  Nothing was omitted from any retained span, so the package carries no omission record.  The retained lines cite 4 works, whose entries are transcribed verbatim from the article's own bibliography source in the e-print and delivered with short editorial labels recorded in the item's reference mapping.  No retained line contains a figure environment, an image-inclusion command or drawing code, so no image is delivered and the package declares no asset dependency.  Editorial formatting only, in addition: an AMS \texttt{amsart} preamble that restates the article's own declarations and macros used by the retained lines, namely the environments \texttt{prop}, \texttt{lem}, \texttt{thm} on separate counters, together with the standard packages those lines need.  The retained environments print in this document as Proposition 1, Lemma 1, Theorem 1 in order of appearance, whereas the article prints the same items with the numbering of its own full document.  The complete native source of the article as distributed in the arXiv e-print for this exact version is bundled unchanged as \texttt{sources/195368/source0.tex}.
\begin{prop}\label{initial}
Let $Q$ be a $K$-quaternion  algebra and $\vf$ a quadratic form over $K$. Let $\pi=\N_Q$.
Let  $(A,\s)= \Ad(\vf)\otimes (Q,\can_Q)$. Then $\G(A,\s)=\G(\vf\otimes \pi)$ and $\Hyp(A,\s)=\Hyp(\vf\otimes \pi)$. Furthermore, for any field extension $K'/K$, we have $${\bf PSim}(A,\s)(K')/R\simeq {\bf PSim}(\vf\otimes\pi)(K')/R\,.$$
\end{prop}

    \begin{lem}\label{24}
      Let $\pi$ a $2$-fold Pfister form. Let $\psi$ be a $6$-dimension quadratic form and $\vf=\pi\otimes \psi$ in $\I^4K$ over $K$. 
      Then, for every $c\in \G(\vf)$, there exists a field extension $M/K$ which is trivial, quadratic or biquadratic such that $\vf_{M}$ is hyperbolic and $c\in\sq{K}\cdot \N_{M/K}^\ast$. In particular, $\G(\vf)=\sq{K}\cdot\Hyp(\vf)$.
    \end{lem}

\begin{thm}\label{6}
    Let $(A,\s)$ be a central simple algebra of $\deg(A)=12$ with symplectic involution. Assume that $\ind(A)\leq 2$ and trivial discriminant. Then  ${\bf PSim}(A,\s)$ is $R$-trivial.
\end{thm}

\begin{proof}
     As $(A,\s)$ be a central simple algebra with symplectic involution with $\ind(A)\leq 2$, by \cite[Prop. 3.4]{BU18} we have $(A,\s)\simeq \Ad(\vf)\otimes (Q,\can_Q)$ for some $6$-dimensional quadratic form $\vf$ over $K$ and $Q$ a quaternion algebra over $K$. 
     Write $Q=(a,b)$, with $a,b\in\mg{K}$. Then its norm form is $\lla a,b\rra$.
 
Let $\psi:=\lla a, b \rra \otimes \vf$.
Denote $\disc(\vf)=c\sq{K}$. 
Then $\vf \equiv \lla c \rra \mod \I^2K$ and hence $\psi=\lla a, b\rra \otimes \vf \equiv \lla a, b ,c\rra \mod \I^4K$. 
By \cite[Main Theorem]{ST04}, we have $\Delta(\s)=[Q]\cup(\disc(\vf))$ in $H^3(K,\mu_2)$.
By our hypothesis, $\Delta(\s)=0$.
Hence 
$(a)\cup(b)\cup (c)=(\disc(\vf))\cup[Q]=0$.   
Hence, by \cite[Theorem 5.1]{MS91} together with \cite[Theorem 4.1]{Mil70}, we have that $\lla a,b,c\rra$ is hyperbolic.

 
 
Hence, $\psi\in \I^4K$. 
By \Cref{24}, it follows that $\G(\psi)=\sq{K}\cdot\Hyp(\psi)$. 
By \Cref{initial}, we have $\G(A,\s)=\G(\psi)$ and $\Hyp(A,\s)=\Hyp(\psi)$. 

 For any field extension $K'/K$, if $c\in \G(\psi_{K'})$ by \Cref{24}, then there exists a field extension $M/K'$ which is trivial, quadratic or biquadratic such that $\psi_M$ is hyperbolic and $c\in\sq{K'}\cdot \N_{M/K'}^\ast$. In particular, $\G(\psi_{K'})=\sq{K'}\cdot\Hyp(\psi_{K'})$ and hence ${\bf PSim}(\psi)$ is $R$-trivial. Thus ${\bf PSim}(A,\s)$ is $R$-trivial.
\end{proof}

\begin{thebibliography}{99}
\bibitem[BU18]{BU18} K.J.~Becher. T.~Unger. Weakly hyperbolic involutions.\emph{Expositiones Math.} 36 (2018), 78--97.
\bibitem[MS91]{MS91} A.S.~Merkurjev. A.A.~Suslin. Norm residue homomorphism of degree three.(Russian) \emph{Izv. Akad. Nauk SSSR Ser. Mat.} 54 (1990), 339--356; translation in Math. USSR-Izv. 36 (1991), 349--367
\bibitem[Mil70]{Mil70} J.~Milnor. Algebraic $K$-theory and quadratic forms. \emph{Invent. Math.} 9 (1969/70), 318–344.
\bibitem[ST04]{ST04} A.~Serhir, J.P.~Tignol. The discriminant of a decomposable symplectic involution. \emph{J. Algebra} 273 (2004), 601--607.
\end{thebibliography}
\end{document}
